\section{Design of \algo}
\label{sec:algorithm_design}
In this section, we first discuss our motivation towards the design of \algo in Section~\ref{subsec:motivation}. We then provide a brief summary of the general workflow of \algo in Section~\ref{subsec:workflow}. After that, we discuss the three key designs of \algo to self-adjust its design components in response to dynamically detected workload characteristics of the data stream, including Regular Stream Characteristics Detection in Section~\ref{subsec:auto-selection}, Automatic Design Choice Selection in Section~\ref{subsec:regular-detection}, Flexible Algorithm Migration in Section~\ref{subsec:migration}.
% This section elucidates the comprehensive algorithmic design of \algo. The algorithm is engineered for flexible adaptability, allowing it to self-adjust its design components in response to dynamically detected workload characteristics of the data stream and user-defined optimization objectives. This adaptability enables \algo to capitalize on the strengths of various design choices, ensuring enhanced and stable clustering performance tailored to specific optimization goals.

\subsection{Motivation} 
\label{subsec:motivation} 
Despite numerous \dsc algorithms, all commonly incorporate the four key design components mentioned earlier. As summarized in Table~\ref{tab:design_comparison}, Wang et al.~\cite{wang2023sesame} revealed that no single design choice consistently guarantees good performance across varying workload characteristics and optimization targets. Based on these findings, Wang et al. proposed \benne, the first \dsc algorithm with a modular architecture that selects different design options from four design aspects. While \benne achieved state-of-the-art performance, it struggles with changes in stream characteristics where a fixed modular composition may not be effective. For instance, when the frequency of cluster evolution of the data stream suddenly increases, using the old configuration designed for low cluster evolution frequency can result in poor performance. To address this issue, we designed \algo, a modular and self-optimizing data stream clustering algorithm. Unlike \benne, \algo dynamically detects changes in workload characteristics and reconfigures itself with the best modular composition.

\subsection{General Workflow}
\label{subsec:workflow}
The \algo's adaptability is rooted in the flexible configuration of four core components: the summarizing data structure (\textit{struc.}), window model (\textit{win.}), outlier detection mechanism (\textit{out.}), and refinement strategy (\textit{ref.}). \algo consists of four main functions to integrate these core components together. Specifically,
\emph{Insert Fun.(...)} adds the new streaming data to the current summarizing data structure, \emph{Window Fun.(...)} manages the data points in the summarizing data structure, \emph{Outlier Fun.(...)} assesses if the current input point qualifies as an outlier. \emph{Refine Fun.(...)} runs the offline clustering before obtaining the final results.

% \emph{(2) Window Function:} 
% The \emph{Window Fun.} (executed at Line 11 in Algorithm~\ref{alg:algo}) manages the data points in the summarizing data structure ($s$), in accordance with the selected window model ($win.$). This function ensures that the algorithm's focus remains on the most recent and relevant data points, thereby allowing it to adapt to changes in the data distribution effectively.
\begin{comment}
\emph{(3) Outlier Detection:}
When the outlier detection mechanism ($out.$) is activated, the algorithm invokes the \texttt{Outlier Fun.} to assess if the current input point ($p$) qualifies as an outlier. This function is executed between Lines 12-15 in Algorithm~\ref{alg:algo}. If the point is determined not to be an outlier, it is subsequently inserted into the summarizing data structure ($s$) and the structure is updated accordingly (executed at Line 15). Conversely, if the outlier detection mechanism is deactivated, the input point is directly inserted into the summarizing data structure, as indicated at Line 17 in Algorithm~\ref{alg:algo}.

\emph{(4) Incremental Refinement:}
Should the refinement strategy ($ref.$) be configured to \incrementalRefinement, the algorithm invokes the \emph{Refine Fun.} each time new data points are processed. This is executed between Lines 18-19 in Algorithm~\ref{alg:algo}. This strategy ensures that the clustering model is continuously updated and remains responsive to dynamic changes in the data distribution. However, this approach may necessitate more frequent updates, thereby potentially increasing computational overhead.

\emph{(5) One-shot Refinement:}
Upon completion of the input data stream processing, the algorithm transitions to the offline phase. If the refinement strategy ($ref.$) is configured as \oneshotRefinement, the \emph{Refine Fun.} is executed in a single pass at the conclusion of the stream processing, as indicated between Lines 20-21 in Algorithm~\ref{alg:algo}. This strategy minimizes computational overhead by limiting the frequency of refinement operations. However, it may compromise the adaptability of the clustering model if it fails to capture timely changes in the data distribution.
\end{comment}

Guided by real-time stream characteristics detected from a sample queue (\textit{queue}), predefined threshold values (\textit{thresholds}), and a performance objective (\textit{o}), \algo dynamically tailors its components to meet the specialized requirements of diverse applications and data streams. In particular, the Regular-detection Fun (Section~\ref{subsec:regular-detection}) detects the changes of characteristics in the current data stream, Auto-selection Fun (Section~\ref{subsec:auto-selection}) selects the best modular configurations based on the detected changes of characteristics, Flexible-migration Fun (Section~\ref{subsec:migration}) migrates the clustering results stored in the old modular configuration to the new one. 

\begin{algorithm}[t]
\scriptsize
\caption{Execution flow of \algo.}
\label{alg:algo}
    \KwData{$p$ \tcp{Input point}}
    \KwData{$c$ \tcp{Temporal Clusters}}
    \KwData{$struc.$ \tcp{Type of summarizing data structure}}   
    \KwData{$win.$ \tcp{Type of window model}}
    \KwData{$out.$ \tcp{Type of outlier detection mechanism}}
    \KwData{$ref.$ \tcp{Type of refinement strategy}}
    \KwData{$queue$ \tcp{Batch of stream for detection}}
    \KwData{$thresholds$ \tcp{Identifying characteristics' changes}}
    \KwData{$o$ \tcp{User's primary performance objective}}
    \tcp{Online Phase}
    Initializing parameters and design choices \;
    \While{!stop processing of input streams}{
      \eIf{\text{$queue$ is full}}{
            $struc\_{old}.$ = $struc.$ \;
            $characteristics$ = \emph{\textcolor{blue}{Regular-detection Fun.(...)}} \;
            $struc.$, $win.$, $out.$, $ref.$ = \emph{\textcolor{blue}{Auto-selection Fun.(...)}} \;
            \emph{\textcolor{blue}{Flexible-migration Fun.(...)}}\;
            Empty $queue$ \;
        }{
            $queue$.push($p$) \;
        } 
        \emph{\textcolor{blue}{Window Fun. (...)}}\;
        \eIf{\text{$out.$ != \nod}}{
            $b$ $\gets$ \emph{\textcolor{blue}{Outlier Fun. (...)}}\;
            \If{\text{$b$ = false}}{
                \emph{\textcolor{blue}{Insert Fun. (...)}}\tcp{Insert $p$ to $s$ and update $s$}            
            }
        }{
            \emph{\textcolor{blue}{Insert Fun. (...)}}\tcp{Insert $p$ to $s$ and update $s$}
        }
        \If{\text{$ref.$ = \incrementalRefinement}}{
            \emph{\textcolor{blue}{Refine Fun. ($ref.$)}};
        }
    }
    \tcp{Offline Phase}
    \If{\text{$ref.$ = \oneshotRefinement}}{
        \emph{\textcolor{blue}{Refine Fun. ($ref.$)}};
    }
\end{algorithm}

Algorithm~\ref{alg:algo} outlines the high-level execution flow of \algo. To leverage advantageous design options under different workload characteristics and optimization targets, \algo operates using a bifurcated execution strategy, comprising an online phase and an offline phase.
In the online phase, the algorithm sequentially processes incoming data points. Initially, parameters and design choices for the summarizing data structure (\textit{struc.}), window model (\textit{win.}), outlier detection mechanism (\textit{out.}), and refinement strategy (\textit{ref.}) are initialized. For each incoming data point \( p \), it is pushed into the sample queue (\textit{queue}). When the queue is full, the algorithm performs regular detection to assess the stream characteristics. The current design choices are stored as \textit{struc\_old.}, and the characteristics are detected using the \emph{Regular-detection Function}. Based on the detected characteristics, the \emph{Auto-selection Function} dynamically selects the optimal design choices for \textit{struc.}, \textit{win.}, \textit{out.}, and \textit{ref.}. The \emph{Flexible-migration Function} then facilitates the transition to the new design choices, and the queue is emptied after the selection process. The \emph{Window Function} processes the data points within the current window model. If an outlier detection mechanism is specified (\textit{out.} $\neq$ \nod), the algorithm uses the \emph{Outlier Function} to check if the point \( p \) is an outlier. If \( p \) is not an outlier, it is inserted into the summarizing data structure and the structure is updated using the \emph{Insert Function}. If no outlier detection is specified, the point \( p \) is directly inserted into the summarizing data structure. If the refinement strategy is set to \textit{Incremental Refinement}, the algorithm updates the clustering model using the \emph{Refine Function}. In the offline phase, after processing the entire data stream, if the refinement strategy is set to \textit{One-shot Refinement}, the clustering model is updated using the \emph{Refine Function} to finalize the clustering results.

% \emph{(1) Automatic Design Choices Adaptation:} Given the dynamic characteristics of data streams, \algo is designed to adapt its configuration to optimize performance continually. The algorithm accumulates incoming data points in a queue (denoted as $queue$ in Line 3 of Algorithm~\ref{alg:algo}). It then performs real-time analysis of the stream's characteristics using the \emph{Regular-detection Fun.} (executed at Line 5). Based on this analysis, \algo dynamically selects the most suitable components for the current data stream using the \emph{Auto-selection Fun.} (executed at Line 6). Detailed discussions of these functions are available in Sections~\ref{subsec:regular-detection} and~\ref{subsec:auto-selection}. If a change in the summarizing data structure is warranted, \algo transfers the temporal clusters from the old structure ($struc.\_old$) to the new one ($struc.$) using the \emph{Flexible-migration Fun.} (executed at Line 7), as elaborated in Section~\ref{subsec:migration}.

% \emph{(2) Window Function:} 
% The \emph{Window Fun.} (executed at Line 11 in Algorithm~\ref{alg:algo}) manages the data points in the summarizing data structure ($s$), in accordance with the selected window model ($win.$). This function ensures that the algorithm's focus remains on the most recent and relevant data points, thereby allowing it to adapt to changes in the data distribution effectively.

% \emph{(3) Outlier Detection:}
% When the outlier detection mechanism ($out.$) is activated, the algorithm invokes the \texttt{Outlier Fun.} to assess if the current input point ($p$) qualifies as an outlier. This function is executed between Lines 12-15 in Algorithm~\ref{alg:algo}. If the point is determined not to be an outlier, it is subsequently inserted into the summarizing data structure ($s$) and the structure is updated accordingly (executed at Line 15). Conversely, if the outlier detection mechanism is deactivated, the input point is directly inserted into the summarizing data structure, as indicated at Line 17 in Algorithm~\ref{alg:algo}.

% \emph{(4) Incremental Refinement:}
% Should the refinement strategy ($ref.$) be configured to \incrementalRefinement, the algorithm invokes the \emph{Refine Fun.} each time new data points are processed. This is executed between Lines 18-19 in Algorithm~\ref{alg:algo}. This strategy ensures that the clustering model is continuously updated and remains responsive to dynamic changes in the data distribution. However, this approach may necessitate more frequent updates, thereby potentially increasing computational overhead.

% \emph{(5) One-shot Refinement:}
% Upon completion of the input data stream processing, the algorithm transitions to the offline phase. If the refinement strategy ($ref.$) is configured as \oneshotRefinement, the \emph{Refine Fun.} is executed in a single pass at the conclusion of the stream processing, as indicated between Lines 20-21 in Algorithm~\ref{alg:algo}. This strategy minimizes computational overhead by limiting the frequency of refinement operations. However, it may compromise the adaptability of the clustering model if it fails to capture timely changes in the data distribution.

\subsection{Regular Stream Characteristics Detection}
\label{subsec:regular-detection}
Various stream characteristics such as data dimensionality, cluster evolution, and the number of outliers are subject to change in evolving data streams. To make informed decisions about the most appropriate design choices for real-time clustering, \algo must first ascertain the current characteristics of the data stream. Algorithm~\ref{algo:auto_detection} delineates the procedure \algo employs to automatically detect key attributes like dimensionality, cluster evolution, and the number of outliers in the current data stream.

\begin{algorithm}[t]
\scriptsize
\caption{Regular-detection Fun. of \algo}
\label{algo:auto_detection}
\SetKwFunction{FMain}{Regular-detection Fun.}
\SetKwProg{Fn}{Function}{:}{}
\tcc{$queue$: a queue of new input stream data stored for detecting stream characteristic changes, $centers$: previous clustering centers of the stream data, $thresholds$: threshold values for identifying the characteristic}
\Fn{\FMain{$queue, centers, threshold$}}{
    \tcc{Initialize stream characteristics}
    \text{$characteristics$ = null}\;
    \text{$high\_dim\_data$ = $var$ = $num\_of\_outliers$= 0}\;
    \text{$new\_center$ = mean($queue$)}\;
    \For{$s \in queue$}{
        \If{$s.dim  > T_d$}{
            $high\_dim\_data$ += 1 \;
        }
        $var$ = UpdateVariance($var, s, new\_centers$)\;
        \If{$min(dist(s, centers))  > thresholds.dist$}{
            $num\_of\_outliers$ += 1 \;
        }
    }
    \eIf{$high\_dim\_data > num(samples) / 2$}{
        $characteristics.high\_dimension = true$ \;
    }{
        $characteristics.high\_dimension = false$ \;
    }
    \eIf{$var > thresholds.variance$}{
        $characteristics.frequent\_evolution = true$ \;
    }{
        $characteristics.frequent\_evolution = false$ \;
    }
    \eIf{$num\_of\_outliers > num(samples) / 2$}{
        $characteristics.many\_outliers = true$ \;
    }{
        $characteristics.many\_outliers = false$ \;
    }
    Return $characteristics$; 
}
\end{algorithm}

Specifically, the algorithm initializes various counters and variables to store stream characteristics at Lines 3-5 of Algorithm~\ref{algo:auto_detection}. For each new data point in the $queue$ (Line 6), \algo evaluates its dimensionality. If the dimension exceeds the threshold $T_d$, the counter $high\_dim\_data$ is incremented  (Line 7). The variance of the data stream is updated at Line 8. The algorithm also checks whether each data point is an outlier based on its distance to the current clustering centers (Line 9-10). 
After processing all the data points in the $queue$, \algo sets the $characteristics.high\_dimension$ attribute to ``true'' or ``false'' based on the value of $high\_dim\_data$ (Lines 11-14). Similarly, the $characteristics.frequent\_evolution$ attribute is determined based on the calculated variance (Lines 15-19), and the $characteristics.many\_outliers$ attribute is set based on the number of outliers detected (Lines 20-23).

\begin{algorithm}[t]
\scriptsize
\caption{Auto-selection Fun. of \algo}
\label{algo:auto_selection}
\SetKwFunction{FMain}{Auto-Selection Fun.}
\SetKwProg{Fn}{Function}{:}{}
\tcc{$objective$: User's primary performance objective (Accuracy, Efficiency, or Balance); $data$: Input data stream; $characteristics$: Stream characteristics detected by regular-detect Fun.}
\Fn{\FMain{$objective$, $data$, $characteristics$}}{
    \tcc{Initialize  four modular design components}
    \text{$struc.$= $win.$ = $out.$= $ref.$ = null}\;
    \eIf{$objective$ = Accuracy}{
        \tcc{Select modules with high accuracy based on input data stream characteristics}
        \eIf{$characteristics.frequent\_evolution = true$}{
            $struc.$ = \MCs \;
        }{
            $struc.$ = \CFT \;
        }
        \eIf{$characteristics.many\_outliers = true$}{
            $win.$ = \landmark \;
            $out.$ = \buffertimer \;
        }{
            $win.$ = \damped \;
            \eIf{$characteristics.high\_dimension = true$}{
                $out.$ = \buffer \;
            }{
                $out.$ = \buffertimer \;
            }
        }
        $ref.$ = \incrementalRefinement \;
    }{
        \tcc{Choose a structure with high efficiency based on input data stream characteristics}
        \eIf{$characteristics.frequent\_evolution = true$}{
            $struc.$ = \DPT \;
            $win.$ = \landmark \;
        }{
            $struc.$ = \Grids \;
            $win.$ = \sliding \;
        }
        $out.$ = \nod \;
        $ref.$ = \norefinement \;
    }
    
    Return $struc.$, $win.$, $out.$, $ref.$; 
}
\end{algorithm}

\subsection{Automatic Design Choice Selection}
\label{subsec:auto-selection}
Upon receiving the stream characteristics from the \emph{Regular-detection Function}, \algo proceeds to select the most appropriate design choices based on these characteristics and the user-defined optimization objective. The detailed selection process is delineated in Algorithm~\ref{algo:auto_selection}.

If the optimization objective is set to Accuracy (Line 3), \algo selects the summarizing data structure based on the frequency of cluster evolution. Specifically, \MCs is chosen if frequent cluster evolution is detected (Line 6), otherwise \CFT is selected (Line 8). For the window model and outlier detection mechanisms, \algo selects \landmark and \buffertimer if many outliers are present (Lines 10-11). If outliers are scarce, \damped is chosen as the window model, and the choice between \buffer and \buffertimer for outlier detection is influenced by the data's dimensionality (Lines 13-16). The refinement strategy is set to \incrementalRefinement (Line 18). If the optimization objective is set to Efficiency (Line 20), \algo selects either \DPT or \Grids as the summarizing data structure based on the frequency of cluster evolution (Lines 22-25). The window model is also selected accordingly, with \landmark chosen for frequent cluster evolution and \sliding otherwise (Lines 24-25). The outlier detection mechanism is set to \nod (Line 26), and the refinement strategy is set to \norefinement (Line 27).

% The subsection does not cover the case when the user aims to balance between clustering accuracy and efficiency, although the algorithm code does not include this case either. Therefore, the description is consistent with the algorithm in this regard.
 
\begin{algorithm}[t]
\scriptsize
\caption{Flexible-migration Fun. of \algo.}
\label{alg:flexible_migration}
\SetKwFunction{FMain}{Flexible-migration Fun.}
\SetKwProg{Fn}{Function}{:}{}
\tcc{$objective$: User's primary performance objective (Accuracy or Efficiency); $struc.$: New selected summarizing data structure; $struc\_{old}.$: Previous selected summarizing data structure; $out.$: Previous selected outlier detection mechanism}
\Fn{\FMain{$objective$, $struc.$, $struc\_{old}.$}}{
    \If{\text{$struc.$ != $struc\_{old}.$}}{
        Extract the clustering centers $c$ from $struc.$\;
        \eIf{$objective$ = Accuracy}{
            \tcc{Transfer the old clustering results to the new summarizing data structures for the accuracy or balance objective}
            \If{$out.$ = \buffer or \buffertimer}{
                Extract the outliers $o$ from $out.$\;
            }
            Initialize the new summarizing data structure $struc.$ with $c$ and $o$\;
        }{
            \tcc{Sink the old clustering results for the efficiency objective}
            Sink $c$ into the output\;
            Create a blank object for initializing  $struc.$\; 
        }
    }
}
\end{algorithm}
\subsection{Flexible Algorithm Migration}
\label{subsec:migration}
Due to the significant structural differences among various summarizing data structures, it is not feasible to directly transfer clustering information from the old summarizing data structure to the new one. To address this, \algo employs a migration function, delineated in Algorithm~\ref{alg:flexible_migration}. 
% This approach ensures that \algo can smoothly transition between different summarizing data structures while optimizing for the user-defined objective.
Specifically, if the newly selected summarizing data structure (\textit{struc.}) differs from the previously selected one (\textit{struc\_old.}) (Line 2), the algorithm adapts as follows:
For the \textbf{Accuracy Objective} (Lines 4-7), \algo extracts the clustering centers (\textit{c}) from the old summarizing data structure (Line 3). Additionally, if the old outlier detection mechanism (\textit{out.}) is either \buffer or \buffertimer, outliers (\textit{o}) are also extracted (Line 6). These centers and outliers are then used to initialize the new summarizing data structure (\textit{struc.}) (Line 7).
For the \textbf{Efficiency Objective} (Lines 8-10), \algo extracts the clustering centers (\textit{c}) from the old summarizing data structure and sinks them into the output to avoid computational overhead (Line 9). A new, empty object is created to initialize the new summarizing data structure (\textit{struc.}) (Line 10).


\subsection{Further Implementation Details} 
\algo is designed as a three-threaded pipeline to process data streams, approximating a realistic computational environment. Inter-thread communication is facilitated through a shared-memory queue, reducing latency associated with network transmissions. The \emph{Data Producer Thread} loads benchmark workloads into memory and enqueues each data point into a shared queue. To simulate high-throughput, the input arrival rate is immediate, eliminating idle time for the algorithm. The \emph{Data Consumer Thread} executes the \dsc algorithm, processing the data stream by dequeuing input tuples from the shared queue and generating temporal clustering results. Efficiency metrics are captured in this thread for consistent comparison of various \dsc algorithms. The \emph{Result Collector Thread} stores the temporal clustering results from the Data Consumer Thread. It computes accuracy metrics to minimize interference with efficiency measurements. Clustering quality is evaluated using purity~\cite{10.1145/1557019.1557115}, and the ability to handle \cef is assessed using CMM~\cite{CMM:2011}.